\begin{itemize}
\item \textbf{How many workers are part of the garbage collection and cleaning service in Olot?}
Olot's garbage collection and cleaning service has a fixed workforce of 38 workers, distributed in different shifts and routes to ensure coverage of the entire city.
\item \textbf{How many shifts is the working day divided into?}
The working day is divided into 3 shifts: morning, afternoon, and night. The shifts are structured to meet operational needs, with most of the work done in the morning and at night.
\item \textbf{How many waste collection routes are established?}
There are currently 6 established routes: 2 operate in the morning and 4 at night, covering different areas of the city to ensure efficient collection.
\item \textbf{How many workers are assigned per route?}
Each route typically has 3 workers assigned. However, on some morning routes, only 2 workers may be assigned depending on the workload.
\item \textbf{What is the general condition of the collection trucks?}
The trucks are over 10 years old and have suffered considerable wear due to intensive use. They often experience issues such as corrosion, mechanical wear, and hydraulic system failures.
\item \textbf{How many liters of diesel are consumed monthly?}
The collection service consumes approximately 10,000 liters of diesel per month to power the fleet of trucks, machines, and cars.
\item \textbf{What is the percentage of men and women working in the service?}
Currently, the service has 24% women and 76% men in its workforce.
\item \textbf{Who is responsible for managing waste collection and cleaning in Olot?}
The service is managed by the private company IFGA-URBASER, operating under contract with the Olot City Council for waste collection and street cleaning.
\item \textbf{What type of collection system is used? Is it door-to-door or another model?}
Olot uses a centralized collection system with containers distributed throughout the city. However, in the urban core, a door-to-door collection system is implemented to improve waste management in higher-density areas.
\item \textbf{How many tons of residual waste are collected monthly?}
Each month, between 700 and 750 tons of residual waste are collected. This amount decreases slightly in summer due to seasonal consumption variations.
\item \textbf{How many tons of packaging waste are collected monthly?}
Each month, between 130 and 140 tons of packaging waste are collected. As with other waste fractions, packaging waste generation is lower in summer.
\item \textbf{How many tons of organic waste are collected monthly?}
Each month, between 230 and 240 tons of organic waste are collected. This amount varies depending on the season and consumption habits of the population.
\item \textbf{On which days of the week do the workers operate?}
Workers operate every day of the week. However, on Sunday mornings, only 4 workers are on duty for essential services.
\item \textbf{What are the most common issues with the collection trucks?}
The most common issues include broken lights, corrosion on the truck body, and failures in hydraulic systems, which can cause service delays.
\item \textbf{How many trips are made to the landfill each week?}
Between 15 and 20 trips to the landfill are made weekly, depending on the volume of waste collected and the capacity of the trucks each day.
\item \textbf{Which areas in Olot generate the most waste?}
The areas with the highest waste generation are the city center and the Sant Miquel neighborhood, especially on market days and during events.
\item \textbf{What types of waste containers exist in Olot, and how are they distributed?}
Olot has waste, packaging, and organic matter containers at each collection point. Additionally, strategic points include cardboard and glass containers to facilitate recycling.
\item \textbf{How is the collection of bulky waste and old furniture managed in Olot?}
The collection of bulky waste and old furniture is not managed directly by the City Council but by a public company in the region. Citizens must request the service or take their items to the municipal recycling center.
\item \textbf{What is the recycling rate in Olot?}
The recycling rate in Olot is around 60%. However, there are significant differences between neighborhoods depending on residents' recycling habits.
\item \textbf{Are there penalties for not recycling correctly in Olot?}
Yes, the City Council imposes penalties for non-compliance with waste management regulations. Fines can reach up to €600 in cases such as abandoning large items in public spaces.
\item \textbf{How is street cleaning carried out in Olot, and how frequently?}
Street cleaning is carried out daily using mechanical sweepers. However, on Sundays, this service is not performed regularly.
\item \textbf{Which company is responsible for maintaining and repairing the collection trucks?}
The basic maintenance of the trucks is performed internally by a mechanic and an assistant. However, major repairs are outsourced to specialized workshops due to a lack of in-house equipment.
\item \textbf{How does tourism affect waste generation in Olot?}
Tourism has a moderate impact on waste generation. However, during weekends and peak tourist periods, there is a noticeable increase in packaging and organic waste in tourist areas and restaurants.
\item \textbf{What campaigns have been implemented in Olot to improve selective waste collection?}
Several awareness campaigns have been carried out, including informational campaigns, inspections, and penalties to improve waste separation. However, their overall impact has been limited.
\end{itemize}